\section{Introduction and correction target}

Ishikawa and Shibata (2021) study R\&D competition and cooperation in a two-firm oligopoly with asymmetric R\&D spillovers. Their framework combines a linear-product-market stage with a first-stage R\&D choice, and it allows firms' objectives to place weight on relative profit through a parameter $\lambda$. The paper uses the resulting closed-form equilibria to compare noncooperative and cooperative R\&D across spillover configurations and degrees of market competitiveness. A central part of the numerical and interpretive analysis is therefore not the existence of the two R\&D equilibria, but the location of the parameter boundaries at which the ranking between them changes. These boundaries are summary objects: they compress a multidimensional comparative-static problem into statements about which regime obtains for a given spillover configuration. That compression is useful only if the implied equality set is correctly characterized. A threshold can therefore be substantively wrong even when every equilibrium formula substituted into it is correct. This distinction matters especially under asymmetric spillovers, because replacing a two-dimensional zero set by a scalar statistic such as average spillovers can discard the very asymmetry that the model is designed to study.

This note addresses a narrow issue in that boundary analysis. The equilibrium formulas themselves are not the source of the problem. Independent reconstruction from the primitives reproduces the published output equilibrium and the closed-form noncooperative and cooperative R\&D solutions used below. The correction concerns the threshold characterizations derived from those formulas: the claim that an individual competition--cooperation trigger is independent of the R\&D-cost parameter $\gamma$, the representation of the aggregate equality boundary by a scalar threshold in average spillovers, and the fixed-cutoff classification that follows from that representation. These objects correspond to the discussion surrounding Eq. (38), Observation 1, Eqs. (43)--(46), and Observations 4--5 of Ishikawa and Shibata (2021). The distinction between an equilibrium error and a threshold error is also important for the appropriate remedy. If the closed-form R\&D equilibria were wrong, the model would have to be solved again before any downstream result could be trusted. Here the opposite is true: the published solutions provide the inputs for the correction. The task is to evaluate their equality sets without imposing the additional invariance and scalar-threshold restrictions used in the numerical interpretation. This makes the correction narrow, but it also makes it directly checkable from the original formulas.

We establish three results. First, the individual equality boundary is not $\gamma$-invariant in general. On the asymmetric slice $\lambda=0$ and $\beta_2=0$, the cleared equality equation has an exact regular root at $(\beta_1,\gamma)=(1/3,44/27)$, and the implicit-function derivative of the root with respect to $\gamma$ equals $-9/59$. At the same time, the equality condition on the symmetric slice $\beta_1=\beta_2=b$ factors so that the crossing is
\[
 b^*(\lambda)=\frac{2+\lambda-\lambda^2}{\lambda^2-3\lambda+4},
\]
which is exactly independent of $\gamma$. This special cancellation explains why selected symmetric or fixed-point numerical comparisons can suggest invariance even though the boundary is not invariant away from symmetry.

Second, the aggregate equality boundary is not generally the straight line implied by a scalar average-spillover threshold. At $\lambda=0$, write $s=\beta_1+\beta_2$ and $d=\beta_1-\beta_2$. On the published line $s=1$, the cleared aggregate equality polynomial is proportional to $d^2(18\gamma+d^2-9)$. Under the maintained condition $\gamma\geq1$, every regular asymmetric point on that line therefore fails the aggregate equality condition. We separately show that the corrected aggregate zero set itself varies with $\gamma$ at a regular asymmetric point. Thus the problem is not merely that a displayed line is slightly misplaced: the scalar reduction suppresses economically relevant dependence on asymmetry and the R\&D-cost parameter.

Third, the threshold error changes the qualitative classification for admissible parameters. At $(a,c,\gamma,\lambda,\beta_1,\beta_2)=(100,50,10,0,49/50,0)$, the published average-spillover rule places the point on the noncooperative-R\&D side because $(\beta_1+\beta_2)/2=0.49<1/2$. The retained closed forms instead give aggregate noncooperative R\&D strictly below aggregate cooperative R\&D. The sign difference is exact, not a floating-point artifact. Since the relevant rational equilibrium functions are regular and continuous at the point and both inequalities are strict, the disagreement persists on an open neighborhood. The misclassified set therefore has positive measure.

The scope of the correction is correspondingly limited. This limitation is not a weakness of the argument but part of its identification: each corrected statement can be traced to a specific equality set, while results that do not use that equality set need not be reopened. The output-stage solution and the published closed-form R\&D equilibria are retained. The displayed $\gamma=50$ benchmark values in Table 1 can be reproduced; what fails is their interpretation as supporting a general $\gamma$-invariant threshold. Likewise, the correction does not establish a reversal of a social-welfare ranking, and it does not show that the paper's broader directional conclusion concerning market competitiveness is globally reversed. The qualitative pattern reported in Observation 6 also survives recomputation at the published calibration. The rest of the note first defines the retained comparison objects, then gives the corrected individual and aggregate boundaries, and finally states precisely which conclusions change and which survive.
